
Analysis across 20 LLMs confirms tight coupling (r = 0.993--0.998, p < $2\times 10^{-17})$ across reliability (TrustfulQA), robustness (AdvBench), and fairness (BOLD) benchmarks, far stronger than moderate correlations (r > 0.3) hypothesized from shared root causes. Coupling persists across model scales (r > 0.98 within size strata), challenging the independent dimensions assumption underlying current trustworthiness evaluation.

However, clustering analysis failed to validate failure mode taxonomy: hierarchical clustering produced perfectly stable groups (100\% bootstrap consistency) but poor separation (silhouette = 0.274 < 0.5). This establishes a methodological constraint---3-benchmark evaluation cannot resolve distinct failure modes despite adequate statistical power for correlation detection.

Near-perfect correlations admit two competing explanations. The \textbf{Unified Capability Hypothesis} posits that reliability, robustness, and fairness---as currently operationalized---measure the same underlying construct rather than orthogonal dimensions. The \textbf{Insufficient Resolution Hypothesis} posits that distinct failure modes exist but require broader coverage (5--10 benchmarks) to reveal structure invisible in the 3-benchmark subset. These hypotheses make opposing predictions testable via concrete future work:

\begin{enumerate}
\item \textbf{Expand to 5--10 benchmarks} (ToxiGen, BBQ, HellaSwag, GSM8K, HumanEval). Unified Capability predicts r > 0.95 for >80\% of pairs; Insufficient Resolution predicts mean r < 0.85 with $\geq 30$\% of pairs < 0.7.
\end{enumerate}

\begin{enumerate}
\item \textbf{Factor analysis on 10-benchmark data.} If first component explains >90\% variance, unified construct confirmed. If 3+ factors needed for 70\% variance, distinct dimensions emerge.
\end{enumerate}

\begin{enumerate}
\item \textbf{Test alternative clustering metrics} (Calinski-Harabasz, Davies-Bouldin) and relaxed thresholds (silhouette > 0.2 for r > 0.95 data) to determine whether bootstrap-silhouette divergence reflects metric inappropriateness or true absence of separation.
\end{enumerate}

\begin{enumerate}
\item \textbf{Intervention validation on highest-correlation pair.} Test whether calibration training targeting TrustfulQA improves AdvBench given r = 0.998 coupling. If improvement transfers (Cohen's d > 0.5 for both), validates shared root cause mechanism.
\end{enumerate}

Results shift the conversation from ''score each dimension independently'' to ''analyze correlation structure first.'' If coupling generalizes beyond the 3-benchmark sample, dimension-independent evaluation may fundamentally misspecify the problem, and practitioners should focus on interventions improving general model quality rather than dimension-specific fixes. However, interpretation remains uncertain pending broader measurement (5--10 benchmarks) and intervention validation. The observed coupling (r > 0.99) is robust within the sample, but distinguishing unified constructs from insufficient resolution requires experimental follow-up.
